%##############################################################################
%  PREDECLARED ANALYSIS PROTOCOL
%  Cheetah connectivity project -- Liam Crettol, Weber State University
%
%  Build:  latexmk -pdf protocol.tex
%
%  This is a standalone signable document. It also drops straight into the
%  manuscript as Appendix A1: replace the body of A1_protocol.tex with
%  \input{../../protocol/protocol_body.tex} once this is signed, or just paste
%  everything between \begin{document} and \end{document}.
%
%  HOW TO USE IT
%  1. Work through every line marked DECISION REQUIRED. Defaults are supplied
%     with rationale; you can accept them, but accept them deliberately.
%  2. Date and lock section 0, then tag the commit (git tag -a protocol-v1.0).
%  3. Post to OSF for an independent timestamp (free, ~10 minutes). This is the
%     only external verification of the lock date, so do not skip it.
%  4. From that point, changes go in the amendment log in section 19. Never
%     edit a locked value silently.
%##############################################################################

\documentclass[11pt,letterpaper]{article}
\usepackage[margin=1in]{geometry}
\IfFileExists{lmodern.sty}{\usepackage{lmodern}}{}
\usepackage[T1]{fontenc}
\usepackage[utf8]{inputenc}
\usepackage[protrusion=true,expansion=false]{microtype}
\usepackage{booktabs}
\usepackage{longtable}
\usepackage{array}
\usepackage{enumitem}
\usepackage{siunitx}
\sisetup{group-separator={,}}
\usepackage{xcolor}
\usepackage[hidelinks]{hyperref}
\usepackage{titlesec}
\usepackage{fancyhdr}

\definecolor{navy}{HTML}{1F3864}
\definecolor{decide}{HTML}{B04A00}
\titleformat{\section}{\large\bfseries\color{navy}}{\thesection}{0.6em}{}
\titleformat{\subsection}{\normalsize\bfseries}{\thesubsection}{0.6em}{}

\pagestyle{fancy}\fancyhf{}
\lhead{\small Predeclared analysis protocol}
\rhead{\small v0.1 draft \textbullet\ \today}
\cfoot{\small\thepage}
\renewcommand{\headrulewidth}{0.4pt}

% DECISION REQUIRED marker: renders orange so open decisions cannot hide.
\newcommand{\dec}[1]{\par\vspace{2pt}\noindent\textcolor{decide}{\textbf{DECISION REQUIRED:} #1}\par\vspace{2pt}}
\newcommand{\lockd}[1]{\textbf{#1}}
\newcommand{\structconn}{potential structural connectivity}

\begin{document}

\begin{center}
{\LARGE\bfseries\color{navy} Predeclared Analysis Protocol}\\[6pt]
{\large Changing connections: temporal \structconn, unprotected pinch points,\\
and monitoring priorities for free-ranging cheetahs in southern Africa}\\[10pt]
{Liam Crettol \textbullet\ Weber State University}\\[4pt]
{Version 0.1 (draft) \textbullet\ \today}
\end{center}

\vspace{6pt}
\hrule
\vspace{10pt}

\noindent\textbf{What this document is.} A complete statement of the analytical
choices for this project, fixed before any connectivity output is generated or
inspected. Its purpose is to remove researcher degrees of freedom from a study
design that contains a very large number of individually defensible options.
Every value below was chosen without reference to any result.

\noindent\textbf{What it is not.} It is not immutable. Changes are permitted and
expected; they are recorded in the amendment log (\S\ref{sec:amend}) with a date
and a reason, and the log is published with the manuscript. A disclosed
amendment is sound practice. A silent change is not.

\vspace{10pt}
\hrule

%==============================================================================
\section{Document control and sign-off}

\begin{tabular}{@{}p{4.2cm}p{10.5cm}@{}}
\toprule
Protocol version & 0.1 (draft) \\
Date drafted & \today \\
Date locked & \rule{5cm}{0.4pt} \\
Author & Liam Crettol \\
Author signature & \rule{5cm}{0.4pt} \\
Git lock tag & \texttt{protocol-v1.0} \quad commit: \rule{3cm}{0.4pt} \\
Independent timestamp & OSF DOI: \rule{4cm}{0.4pt} \\
Amendments to date & 0 \\
\bottomrule
\end{tabular}

\vspace{8pt}
\noindent The lock recorded above attests that this protocol was fixed before any
connectivity, corridor, current-flow, protected-area overlap, or prioritisation
result had been generated or viewed. Because this is single-author work with no
countersignature, the verifiable evidence of the lock date is external: the
signed git tag and the OSF timestamp, both of which are dated independently of
the author.

%==============================================================================
\section{Purpose and contribution}

This project models how \structconn\ between established cheetah range cores has
changed across southern Africa between 2012 and 2024, identifies which pinch
points remain important across plausible modelling assumptions, characterises
their protected-area status, and derives a ranked set of monitoring priorities.

Two boundaries define the contribution and both are binding.

\begin{enumerate}[leftmargin=*]
\item \textbf{This is not a habitat suitability study.} Regional current and
projected cheetah habitat suitability, including protected-area mismatch, is
already published. Suitability may enter this analysis only as an input to a
resistance scenario, never as a result.
\item \textbf{Range cores are held fixed at the historical baseline.} No
comparable range reassessment exists after the 2010--2016 evidence period.
Cores therefore do not move between snapshots, and all measured temporal change
is landscape-side. This must be stated in the abstract, not only the methods.
\end{enumerate}

%==============================================================================
\section{Objectives and hypotheses}

\subsection{Objectives}
\begin{enumerate}[leftmargin=*,label=O\arabic*.]
\item Quantify change in modelled \structconn\ between fixed range cores across
four snapshots spanning 2012--2024.
\item Identify pinch points that remain important across the full sensitivity
grid, and characterise their protected-area status.
\item Derive and stress-test a ranked set of monitoring priorities.
\end{enumerate}

\subsection{Hypotheses as locked}

\textbf{H1 (revised from the workbook).} Connectivity loss is concentrated on a
minority of network links rather than distributed across the network, and the
affected links are disproportionately those with high betweenness in the
baseline network.

\emph{Note on the revision.} H1 as originally drafted predicted that cells with
larger increases in human pressure would show larger increases in resistance.
That is true by construction, since the pressure layers are the resistance
inputs. It cannot fail and therefore is not a hypothesis. The revised form asks
a question the model can answer negatively: pressure change might fall on
redundant links, in which case network-level connectivity is robust to it and
the conservation implication reverses.

\emph{Refuted if:} connectivity loss is approximately uniform across links, or
loss falls preferentially on low-betweenness links.

\textbf{H2.} High-value connectivity area is disproportionately located outside
formal protected areas.

\emph{Refuted if:} the outside-PA share of high-current area does not exceed the
predeclared null (\S\ref{sec:pa}) at the stated significance level, or the
result does not hold across current-flow thresholds.

\textbf{H3.} A limited subset of transboundary linkages remains important across
all resistance assumptions, while other apparent corridors are strongly
model-sensitive.

\emph{Refuted if:} agreement is uniformly high (in which case the ensemble adds
nothing) or uniformly low (in which case no robust set exists and the paper's
primary output becomes the uncertainty map, per the decision rule in
\S\ref{sec:stop}).

\textbf{H4.} The highest-priority monitoring cells occur where structural
importance coincides with old or sparse cheetah evidence and rapid landscape
change.

\emph{Refuted if:} no stable top tier persists across the predeclared weight
sets.

\lockd{H1 revision confirmed.} The original form was true by construction
(pressure layers are the resistance inputs) and could not fail, so it was not a
hypothesis. The revised form above is retained as stated.

%==============================================================================
\section{Study extent and exclusions}

\textbf{Primary extent.} The confirmed and likely free-ranging cheetah range
envelope across Namibia, Botswana, Zimbabwe and South Africa, as delineated in
the archived 2010--2016 regional range assessment, buffered outward by
\lockd{\SI{100}{\kilo\metre}} to avoid truncating corridors at the envelope edge.

\emph{Rationale for the buffer.} Least-cost and current-flow results are
sensitive to boundary truncation. A buffer wider than the plausible inter-core
path deviation prevents the analysis boundary from becoming an artificial
barrier. Cores are not defined in the buffer; it exists only so paths can route
through it.

\textbf{Exclusions.} Permanent water bodies; cells with no valid data in any
covariate for any snapshot; areas outside the four primary countries except
where inside the buffer.

\textbf{Scope-control fallback.} If total compute for the full sensitivity grid
exceeds \lockd{120 hours} on available hardware, the extent reduces to the
Namibia--Botswana--Zimbabwe transboundary network. This trigger is numeric and
declared in advance so that a scope reduction is a planned contingency rather
than a response to inconvenient results. South Africa would then be reported as
a validation region only.

\lockd{\SI{100}{\kilo\metre} buffer confirmed.} It is the same order of
magnitude as the individual movement-distance value adopted for the PC
dispersal parameter below, which gives the two figures internal consistency,
and it is wide enough that corridor routing is not truncated at the boundary.

\dec{The 120-hour compute trigger remains provisional pending a pilot run on
the development hardware (GitHub Codespace: 4-core, 16GB RAM, 32GB storage).
Time one full sensitivity-grid unit run (one resistance scenario $\times$ one
core definition $\times$ one snapshot at \SI{1}{\kilo\metre}) end to end in
Circuitscape.jl pairwise mode, multiply by 144 runs, and compare the result
against the 120-hour ceiling before timeline block 5. Amend this value with the
measured figure once the pilot is run.}

%==============================================================================
\section{Coordinate system and working resolution}

\textbf{Coordinate reference system.} Africa Albers Equal Area Conic
(ESRI:102022), an equal-area projection appropriate for continental Africa.
Equal-area is required because the analysis reports areal quantities (protected
area shares, corridor areas, tabulated change).

\lockd{Full WKT confirmed below (Esri's standard definition for ESRI:102022).}

\begin{center}
\fbox{\parbox{0.92\textwidth}{\scriptsize\ttfamily\raggedright
PROJCS["Africa\_Albers\_Equal\_Area\_Conic",GEOGCS["WGS 84",DATUM["WGS\_1984",SPHEROID["WGS 84",6378137,298.257223563,AUTHORITY["EPSG","7030"]],AUTHORITY["EPSG","6326"]],PRIMEM["Greenwich",0,AUTHORITY["EPSG","8901"]],UNIT["degree",0.0174532925199433,AUTHORITY["EPSG","9122"]],AUTHORITY["EPSG","4326"]],PROJECTION["Albers\_Conic\_Equal\_Area"],PARAMETER["latitude\_of\_center",0],PARAMETER["longitude\_of\_center",25],PARAMETER["standard\_parallel\_1",20],PARAMETER["standard\_parallel\_2",-23],PARAMETER["false\_easting",0],PARAMETER["false\_northing",0],UNIT["metre",1,AUTHORITY["EPSG","9001"]],AXIS["Easting",EAST],AXIS["Northing",NORTH],AUTHORITY["ESRI","102022"]]
}}
\end{center}

\textbf{Working resolution.} \lockd{\SI{1}{\kilo\metre}}, with the full pipeline
repeated at \lockd{\SI{2}{\kilo\metre}} and \lockd{\SI{5}{\kilo\metre}} as
resolution sensitivity.

\emph{Rationale.} The occurrence baseline is a \SI{10}{\kilo\metre} assessment
grid; environmental covariates range from \SI{30}{\metre} to \SI{1}{\kilo\metre}.
Modelling at \SI{10}{\kilo\metre} would make a ``pinch point'' a
\SI{100}{\kilo\metre\squared} block, which is not an actionable location and not
a pinch point in any meaningful sense. Modelling at \SI{100}{\metre} would imply
locational precision the occurrence evidence cannot support and would raise
compute by two orders of magnitude. \SI{1}{\kilo\metre} is the coarsest
resolution at which a pinch point is a place a survey team could go, and the
finest at which the claim remains honest. The two coarser repetitions
demonstrate that conclusions are not resolution artefacts.

\textbf{Alignment.} All rasters share a single snap raster, origin, cell size,
extent and NoData convention. Aggregation from finer sources uses mean for
continuous variables and majority for categorical variables; both are recorded
per layer.

%==============================================================================
\section{Temporal design}

\textbf{Snapshots.} \lockd{2012, 2016, 2020, 2024}.

\textbf{Nearest-year rule.} Where a covariate is not available for a snapshot
year, the nearest available year is used and the offset in years is recorded.
The temporal alignment table (one row per covariate, one column per snapshot,
each cell giving source year and offset) is a required deliverable and is
published in the supplement.

\textbf{Maximum tolerated offset.} \lockd{$\pm$2 years}. A covariate that cannot
meet this for a given snapshot is treated as static for the entire analysis
rather than being interpolated. Interpolation is prohibited: it would create
apparent change from an assumption rather than an observation.

%==============================================================================
\section{Range cores}

Three core definitions, all derived from the archived confirmed-range data,
applied identically to every snapshot.

\begin{center}
\begin{tabular}{@{}llll@{}}
\toprule
ID & Evidence included & Minimum patch area & Character \\
\midrule
C1 & Confirmed range only & \SI{2000}{\kilo\metre\squared} & Conservative \\
C2 & Confirmed + likely range & \SI{1000}{\kilo\metre\squared} & Moderate (default for headline figures) \\
C3 & Confirmed + likely + possible & \SI{500}{\kilo\metre\squared} & Inclusive \\
\bottomrule
\end{tabular}
\end{center}

\emph{Rationale for the area thresholds.} Cheetah home ranges on southern
African farmland are reported at the order of hundreds to over a thousand square
kilometres. A core smaller than roughly one home range is not a population
source in any useful sense. The three thresholds bracket that magnitude rather
than asserting a precise value.

\lockd{Area thresholds confirmed.} Marker et al. 2008 report a mean cheetah home
range of \SI{1651}{\kilo\metre\squared} ($\pm$\SI{1594}{\kilo\metre\squared}) on
north-central Namibian farmland; Melzheimer et al. 2020 report
\SI{380}{\kilo\metre\squared} for territory holders, up to
\SI{1600}{\kilo\metre\squared} for floaters spanning multiple territories, and
\SI{650}{\kilo\metre\squared} for females. This confirms the ``hundreds to over
a thousand'' order of magnitude above, and the three thresholds bracket it
without adjustment.

\emph{Amendment: evidence category mapping.} The archived Weise et al. 2017
Dryad raster (\texttt{cheetahsouthernafricadistribution.tif}) does not use
confirmed/likely/possible labels. Its raster attribute table instead codes:
Value 1 = Possible Cheetah Presence; Value 2 = Free Range Cheetah Presence
Buffer; Value 3 = Free Range Cheetah Presence; Value 4 = Fenced Cheetah
Presence. Mapped onto the core definitions above: C1 (confirmed range only)
uses Value 3 only; C2 (confirmed + likely) adds Value 2; C3 (confirmed + likely
+ possible) adds Value 1. Value 4 (fenced populations) is excluded from all
three core definitions, consistent with the free-ranging scope of this study.

Cores are generated by region grouping on the thresholded evidence surface, with
a single morphological closing pass to remove single-cell holes. No manual
editing of core geometry is permitted at any stage.

%==============================================================================
\section{Covariate register and temporal treatment}

Every covariate is classified as \textbf{dynamic} (varies across snapshots and
therefore carries temporal signal) or \textbf{static} (identical in every
snapshot). This classification is locked here and cannot change after results
are seen.

\begin{center}
\small
\begin{tabular}{@{}p{4.6cm}p{2.4cm}p{2.0cm}p{5.2cm}@{}}
\toprule
Covariate & Source & Treatment & Note \\
\midrule
Built-up fraction & GHSL & Dynamic & Epochal; apply nearest-year rule \\
Nighttime radiance & VIIRS annual & Dynamic & Mask known gas flares before differencing \\
Percent tree cover & MODIS VCF & Dynamic & Pin one collection version across all years \\
Percent non-tree vegetation & MODIS VCF & Dynamic & As above \\
Percent bare ground & MODIS VCF & Dynamic & As above \\
Land cover (categorical) & MODIS MCD12Q1 & \textbf{Static} & Context and masking only; never differenced \\
Roads & GRIP4 (primary) & \textbf{Static} & OSM for contemporary detail only \\
Livestock density & GLW4 & \textbf{Static} & Single reference year; exposure only \\
Veterinary fences & national / AHEAD & \textbf{Static} & See \S\ref{sec:fences} \\
Distance to water & HydroSHEDS & \textbf{Static} & \\
Slope / ruggedness & SRTM & \textbf{Static} & \\
Protected areas & WDPA & \textbf{Static} & Overlay only, not a resistance input \\
\bottomrule
\end{tabular}
\end{center}

\textbf{Why categorical land cover is static.} Annual categorical land-cover
products reassign classes between adjacent years for reasons unrelated to
landscape change, and the instability concentrates in exactly the semi-arid
mixed grass and shrub classes that dominate this study area. Differencing them
would manufacture change. Temporal vegetation signal is therefore carried
entirely on the continuous VCF fields.

\subsection{Veterinary fences}
\label{sec:fences}

Veterinary cordon fences in Botswana and Namibia are major deliberate barriers
to large mammal movement and must be represented. Acquisition proceeds in this
order: national veterinary department layers; the AHEAD/WCS KAZA fence
assessment materials; OSM \texttt{barrier=fence} where mapped.

\textbf{Contingency.} If no usable vector layer can be obtained by the end of
timeline block 3, fences are handled as an explicit scenario comparison (network
modelled with and without digitised fence lines from the best available source)
and the limitation is stated prominently. Omitting fences silently is not an
option.

\lockd{Acquisition order and contingency trigger confirmed.} National
veterinary department layers rank first as the most authoritative source, the
AHEAD/WCS KAZA fence assessment second, and OSM \texttt{barrier=fence} last as a
crowdsourced fallback. The trigger is the end of timeline block 3 on the Master
Workflow schedule, not a calendar date: if the acquisition sequence above has
not produced a usable vector layer by that point in the workflow, the
scenario-comparison contingency above is invoked automatically.

%==============================================================================
\section{Resistance scenarios}

Six scenarios. RES-01 to RES-05 are as specified in the project workbook's
Resistance Scenario Register, reproduced by reference and treated as locked.
RES-06 is added here.

\textbf{RES-06 (composite pressure).} Resistance driven by a single published
composite human-modification index rather than by separately weighted pressure
layers. Its purpose is to test whether conclusions depend on the analyst's own
weighting of correlated human-pressure variables, which is the collinearity
concern in the cautions register. It uses no parameters chosen by the author.

\textbf{No preferred scenario.} No scenario is designated as correct, primary,
or best. Results are reported as an ensemble. Where a single surface is required
for illustration, RES-04 (balanced evidence) is used and labelled as
illustrative, not preferred.

\textbf{Transformations.} As specified per scenario in the register. Resistance
values are relative ranks, not empirical coefficients, and are described as such
throughout.

\lockd{RES-06 and the RES-04 illustrative designation confirmed.} Both were
fixed before any resistance surface was produced.

%==============================================================================
\section{Connectivity modelling}

\textbf{Effective distance and corridors.} ArcGIS Pro Distance Accumulation and
Optimal Region Connections, cores as sources, per scenario and per snapshot.

\textbf{Current flow.} Circuitscape.jl in \lockd{pairwise} mode, all core pairs,
with cores as sources and grounds. Pairwise rather than omnidirectional because
cores are defined a priori; an omnidirectional formulation answers a different
question and is not used.

\textbf{Normalisation.} Cumulative current is normalised by the number of core
pairs contributing, so that snapshots and core definitions yielding different
core counts remain comparable. The normalisation formula is fixed here and
applied identically everywhere.

\textbf{Corridor definition.} \lockd{Continuous current surfaces summarised by
quantile.} Thresholded corridor polygons are used only for cartography and are
never the basis of a statistical comparison.

\emph{Why this must be fixed now.} The corridor definition sets the denominator
for the protected-area test in \S\ref{sec:pa}. Choosing it after seeing the
overlap result would let the analyst select the definition that produces the
desired protection gap.

\textbf{High-current definition.} The \lockd{top decile} of normalised current
within the analysis extent.

%==============================================================================
\section{Network metrics}

Cores are nodes; effective distances between core pairs are edge weights.
Reported metrics:

\begin{itemize}[leftmargin=*,nosep]
\item Probability of connectivity (PC) and per-node and per-link importance
(dPC), computed in Conefor or an equivalent implementation.
\item Betweenness centrality on the same graph, used for H1.
\end{itemize}

\textbf{Dispersal distance parameter.} PC requires a dispersal probability
function. Median dispersal probability is set at \lockd{\SI{100}{\kilo\metre}},
with sensitivity at \lockd{50, 200 and \SI{400}{\kilo\metre}}. No single value
is reported alone.

\lockd{\SI{100}{\kilo\metre} confirmed as an order-of-magnitude central value.}
No published cheetah natal-dispersal telemetry study exists for this system.
The best available proxy is Melzheimer et al. 2020, who report communication
hubs spaced \SI{22.9}{\kilo\metre} ($\pm$\SI{4.0}{\kilo\metre}) apart (independently corroborated at
approximately \SI{20}{\kilo\metre} in the Serengeti and Maasai Mara), with
floaters visiting an average of $3.3 \pm 1.9$ hubs --- implying individual
movement typically spans on the order of several tens to \SI{100}{\kilo\metre}
within an extended range. This supports \SI{100}{\kilo\metre} as a defensible
order-of-magnitude central value, while acknowledging it is derived from
hub-spacing data rather than a dedicated dispersal study. Sensitivity at 50,
200 and \SI{400}{\kilo\metre} is reported alongside it; no single value is
asserted as definitive.

%==============================================================================
\section{Temporal comparison}

Change is computed between consecutive snapshots and between the first and last
snapshot, using identical cores, scenarios, resolution and normalisation. Both
absolute and proportional change are reported. Change is attributed to covariate
change by construction; no causal language is used.

%==============================================================================
\section{Protected-area analysis}
\label{sec:pa}

\textbf{Source.} A single named WDPA monthly release, recorded here on locking.
\lockd{Release: \rule{3cm}{0.4pt}}

\textbf{Geometry handling.} Valid polygons only. Point records are excluded and
the count of exclusions is reported. Status and designation fields are filtered
to \lockd{designated, inscribed and established}; proposed sites are excluded.

\textbf{Predeclared null.} The observed share of top-decile current area falling
outside protected areas is compared against a null generated by
\lockd{999 area-standardised spatial permutations of the protected-area
configuration within the analysis extent}, preserving total protected area and
patch size distribution.

\textbf{Primary comparison denominator.} \lockd{The modelled network} (the union
of the top quartile of current), not the full analysis extent. The full extent
is reported as a secondary comparison.

\emph{Rationale.} Comparing against the full extent would largely recover the
known fact that protected areas cover a minority of the landscape. Comparing
against the modelled network asks the sharper question: within the area the
model identifies as connectivity-relevant, is the highest-value subset
disproportionately unprotected?

\textbf{Significance threshold.} \lockd{$\alpha = 0.05$}, reported with the
observed effect size, not as a bare $p$-value.

\textbf{Threshold sensitivity.} The comparison is repeated at the top 5\%, 10\%
and 20\% of current. A result that holds only at one threshold is reported as
threshold-dependent.

%==============================================================================
\section{Sensitivity grid and the definition of robustness}
\label{sec:grid}

\textbf{Full grid at \SI{1}{\kilo\metre}:} 6 resistance scenarios $\times$ 3 core
definitions $\times$ 4 snapshots = \lockd{72 runs}.

\textbf{Resolution sensitivity:} the same grid at \SI{2}{\kilo\metre} and
\SI{5}{\kilo\metre} for the \lockd{first and last snapshots only}
(6 $\times$ 3 $\times$ 2 $\times$ 2 = 72 additional runs), giving
\lockd{144 runs total}.

\emph{Rationale for restricting the resolution sensitivity.} Repeating all four
snapshots at all three resolutions triples compute for little added inference.
The first and last snapshots bracket the change being measured, which is what
the resolution check needs to test.

\textbf{Definition of robust.} A cell is robust if it falls in the top decile of
normalised current in \lockd{at least 80\% of runs} within a given snapshot
(that is, $k \geq 15$ of $n = 18$ scenario-core combinations at
\SI{1}{\kilo\metre}).

\emph{Rationale for $k/n = 0.8$.} The threshold is set high because the
scenarios share input layers and will therefore agree more than genuinely
independent models would. A permissive threshold would report shared-input
agreement as biological robustness. This number is fixed before any agreement
map is produced.

\textbf{Reporting rule.} Cells with agreement below the threshold are reported
as \emph{uncertain}, never as priorities, and appear in the monitoring index
rather than in the corridor recommendations.

\lockd{$k/n = 0.8$ confirmed.} This is the single most consequential number in
the protocol and the one most tempting to relax later. If it is relaxed after
seeing the agreement map, it must appear in the amendment log with the pre- and
post-values and the original result reported alongside.

%==============================================================================
\section{Independent plausibility assessment}

\textbf{Route: plausibility check only.} No resistance surface is fitted to the
occurrence data. Consequently there is no fitted model to validate, and the word
``validation'' is prohibited in describing this analysis (see the claim
dictionary, \S\ref{sec:dict}).

\textbf{Evidence used.} The archived Limpopo cheetah tracking dataset
(nine individuals, 2003--2008) and published camera-trap occurrence evidence
from central-eastern Namibia, both external to the core definition.

\textbf{Test.} Observed locations are compared against availability-matched
random points drawn within the same landscape, asking whether observed locations
fall in lower-resistance and higher-current cells more often than expected.
Autocorrelation within tracks is accounted for by \lockd{subsampling to one
location per individual per 24 hours}.

\textbf{Declared limits.} Nine individuals in one landscape twenty years ago
cannot speak to a four-country network in 2024. A positive result is localised
support only. A negative or null result is reported with equal prominence and
does not invalidate the modelling, but does bound the discussion.

%==============================================================================
\section{Monitoring priority index}

\textbf{Reporting unit.} \SI{10}{\kilo\metre} cells, matching the resolution of
the occurrence evidence rather than the modelling resolution.

\textbf{Components}, each normalised to $[0,1]$ by min-max scaling within the
extent: (i) connectivity importance (normalised current and dPC);
(ii) magnitude of connectivity change 2012--2024; (iii) model uncertainty
(inverse of ensemble agreement); (iv) evidence age; (v) evidence sparsity;
(vi) threat exposure (change in human pressure).

\textbf{Predeclared weight sets.}

\begin{center}
\begin{tabular}{@{}lcccccc@{}}
\toprule
Set & Connectivity & Change & Uncertainty & Evid. age & Evid. sparsity & Threat \\
\midrule
W1 (balanced, headline) & 0.25 & 0.20 & 0.15 & 0.15 & 0.15 & 0.10 \\
W2 (connectivity-led) & 0.45 & 0.15 & 0.10 & 0.10 & 0.10 & 0.10 \\
W3 (knowledge-gap-led) & 0.15 & 0.15 & 0.20 & 0.25 & 0.20 & 0.05 \\
\bottomrule
\end{tabular}
\end{center}

\textbf{Stability testing.} Spearman rank correlation between weight sets across
all cells, and Jaccard overlap of the top 20 cells between weight sets. A cell
enters the recommended tier only if it appears in the top 20 under
\lockd{all three} weight sets.

\textbf{Transparency requirement.} All six component scores are published per
cell alongside the composite, so any reader can reweight and obtain their own
ranking. The index is decision support under stated assumptions, not an
objective ordering.

\lockd{The three weight sets are confirmed and fixed.} If a partner
organisation later supplies operational weights, those are added as a fourth
set in an amendment, not substituted for these.

%==============================================================================
\section{Claim dictionary}
\label{sec:dict}

Permitted and prohibited language, binding on the manuscript.

\begin{center}
\small
\begin{tabular}{@{}p{7.1cm}p{7.4cm}@{}}
\toprule
\textbf{Prohibited} & \textbf{Required instead} \\
\midrule
functional corridor; cheetah corridor & area of high modelled \structconn \\
cheetahs move through / use this area & this area is a modelled structural priority warranting field validation \\
validation; validated & independent plausibility assessment; consistent with \\
conflict hotspot & area of high modelled livestock exposure \\
population decline; caused & spatially coincided with; associated with \\
current cheetah range & historical confirmed-range evidence, 2010--2016 \\
absent; no cheetahs & no records; unknown detection and effort \\
the corridor network & the modelled network under scenario X \\
this resistance surface is correct & results were evaluated across a predeclared ensemble \\
\bottomrule
\end{tabular}
\end{center}

Every prohibited term is checked by text search before submission. This check is
a required item in the pre-submission audit.

%==============================================================================
\section{Decision and stopping rules}
\label{sec:stop}

\textbf{If corridor rankings are unstable} (no robust set emerges under
\S\ref{sec:grid}), the paper's primary output becomes the uncertainty-guided
monitoring design, and the corridor map is demoted to a supporting figure. This
is a legitimate result, not a failure, and is declared here so it cannot be
presented as the plan all along or quietly avoided.

\textbf{If the protected-area result is not significant}, H2 is reported as
unsupported and the discussion addresses why, without threshold shopping.

\textbf{If the plausibility assessment is null}, it is reported in full and the
discussion is bounded accordingly.

\textbf{All scenarios are retained and reported.} No scenario is dropped for
producing inconvenient or uninteresting results.

%==============================================================================
\section{Reproducibility and data handling}

\textbf{Scripting.} All geoprocessing is scripted in arcpy. ModelBuilder is used
for exploration only; no result in the manuscript derives from an unscripted
step.

\textbf{Version recording.} ArcGIS Pro, Julia, Circuitscape.jl, R or Python, and
every dataset version and access date are recorded in the data manifest.

\textbf{Redistribution limits.} Raw WDPA data are not redistributed; users
obtain them independently under Protected Planet terms. Occurrence data are
shared only at the resolution the source archive permits. Derived summaries and
all code are published.

\textbf{Archive.} Code, protocol, manifest and derived non-restricted outputs
are deposited with a DOI on completion.

%==============================================================================
\section{Amendment log}
\label{sec:amend}

Every change to a locked value after the sign-off date is recorded here before
the change is made. The log is published with the manuscript.

\begin{center}
\begin{tabular}{@{}p{1.4cm}p{1.6cm}p{3.4cm}p{3.2cm}p{4.4cm}@{}}
\toprule
No. & Date & Section and value & Changed from / to & Reason, and whether any result had been seen \\
\midrule
1 & & & & \\[10pt]
2 & & & & \\[10pt]
3 & & & & \\[10pt]
4 & & & & \\[10pt]
\bottomrule
\end{tabular}
\end{center}

\vspace{6pt}
\noindent The final column is the one that matters. An amendment made before any
relevant result was viewed is routine. An amendment made afterwards is not
disqualifying, but it must be visible, and the pre-amendment result is reported
alongside the post-amendment result.

\end{document}
